\documentclass[11pt]{article}
% Shared print layout. The web's provisional design and MathJax are independent.
\usepackage[utf8]{inputenc}
\usepackage[T1]{fontenc}
\usepackage[spanish,es-nodecimaldot]{babel}
\usepackage{lmodern}
\usepackage{amsmath,amssymb}
\usepackage{graphicx}
\usepackage{geometry}
\geometry{a4paper,left=20mm,right=20mm,top=25mm,bottom=22mm,headheight=16pt,headsep=8mm}
\usepackage{fancyhdr,lastpage}
\usepackage{enumitem}
\usepackage{xcolor}
\usepackage{pgfplots}
\pgfplotsset{compat=1.18}
\usepgfplotslibrary{fillbetween}
\usepackage[hidelinks,pdfencoding=auto]{hyperref}
\hypersetup{pdflang=es-ES,pdfauthor={},pdfsubject={Matemáticas Aplicadas a las Ciencias Sociales II}}
\pagestyle{fancy}
\fancyhf{}
\fancyhead[L]{\small Matemáticas Aplicadas a las Ciencias Sociales II}
\fancyhead[R]{\small 2.\textsuperscript{o} Bachillerato}
\fancyfoot[L]{\footnotesize Límites y derivadas | Edición revisada 1.0}
\fancyfoot[R]{\footnotesize Página \thepage\ de \pageref{LastPage}}
\renewcommand{\headrulewidth}{0.4pt}
\setlength{\parindent}{0pt}
\setlength{\parskip}{0.5em}
\setlist{itemsep=0.3em,topsep=0.3em}
\newcommand{\apartado}[1]{\subsection*{#1}}
\pgfplotsset{
  revision/.style={
    width=0.94\linewidth,height=7.5cm,
    axis lines=middle,grid=major,grid style={gray!18},
    tick label style={font=\small},label style={font=\small},
    xlabel={$x$},ylabel={$y$},samples=160,
    legend style={at={(0.5,-0.16)},anchor=north,draw=none,font=\small},
    clip=true
  }
}

\hypersetup{pdftitle={Límites y derivadas - Examen revisado}}
\begin{document}
\begin{minipage}{0.17\linewidth}\includegraphics[width=0.8\linewidth]{logo-al-andalus.png}\end{minipage}\hfill
\begin{minipage}{0.79\linewidth}{\Large\bfseries Límites y derivadas}\\[0.4em]
Duración: \textbf{55 minutos}. Puntuación: \textbf{10 puntos}.\end{minipage}

Apellidos y nombre: \hrulefill

Fecha: \hrulefill

Responde razonadamente y numera tus páginas. Se permite el uso de calculadora.
\section*{Ejercicio 1 \hfill 3,2 puntos}
\[f(x)=\begin{cases}x^2+ax+2,&x\leq0,\\(x+b)/(x-1),&x>0,\end{cases}\qquad a,b\in\mathbb R.\]
\begin{enumerate}[label=(\alph*)]
\item (2 puntos) \textquestiondown{}Para qué valores de $a$ y $b$ es continua y derivable en $x=0$?
\item (0,6 puntos) Calcula la recta tangente en $x=2$.
\item (0,6 puntos) Calcula la recta normal en $x=2$.
\end{enumerate}
\section*{Ejercicio 2 \hfill 2,4 puntos}
Calcula las derivadas de \textbf{cuatro} de las seis funciones. Cada apartado vale \textbf{0,6 puntos}. Se corregirán \textbf{las cuatro primeras} respuestas que aparezcan, según orden de lectura.
\begin{enumerate}[label=(\alph*)]
\item $\displaystyle f(x)=\left(\frac{3x^2-5}{2x+3}\right)^3$.
\item $\displaystyle f(x)=\sin^2 x$.
\item $\displaystyle f(x)=\frac{x^3-4}{(x^2+4)^2}$.
\item $\displaystyle f(x)=e^{3x^2-x+2}$.
\item $\displaystyle f(x)=\cos(x^2+2x)$.
\item $\displaystyle f(x)=\frac{x^2}{\sqrt{1+2\ln x}}$.
\end{enumerate}
\clearpage
\section*{Ejercicio 3 \hfill 2,4 puntos}
Calcula \textbf{cuatro} de los seis límites. Cada apartado vale \textbf{0,6 puntos}. Se corregirán \textbf{los cuatro primeros} que aparezcan, según orden de lectura.
\begin{enumerate}[label=(\alph*)]
\item $\displaystyle\lim_{x\to0}\frac{1-\sqrt{1-x^2}}{x^2}$.
\item $\displaystyle\lim_{x\to+\infty}(3^x-x^5)$.
\item $\displaystyle\lim_{x\to+\infty}\frac{5x+7}{\sqrt{2x^2+3}}$.
\item $\displaystyle\lim_{x\to+\infty}\left(\frac{x^2-1}{x^2}\right)^x$.
\item $\displaystyle\lim_{x\to4}\left(\frac7{x-4}-\frac{4x+12}{x^2-4x}\right)$.
\item $\displaystyle\lim_{x\to1}\frac1{x^2-3x+2}$.
\end{enumerate}
\section*{Ejercicio 4 \hfill 2 puntos}
La distancia entre un cohete en descenso y su zona de aterrizaje viene dada por
\[D(t)=\frac{100}{t^2+1},\]
donde $D$ se mide en \textbf{miles de kilómetros} y $t$ en horas desde el comienzo de la fase final de descenso.
\begin{enumerate}[label=(\alph*)]
\item (1 punto) \textquestiondown{}A qué distancia está al comenzar esa fase? \textquestiondown{}Y después de dos horas?
\item (1 punto) \textquestiondown{}A qué valor tiende la distancia cuando $t$ tiende a infinito?
\end{enumerate}
\end{document}
